% EDITION: citation-based shortening, 4 October 2026.
% Earlier audits refer to the source edition; this revision has not been independently audited.
\section{P03: Origin and homogeneous linearization}
\label{sec:P03}
Use the canonical demand-modulated field \eqref{eq:canonical-dm},
with $n\ge5$, $1\le q\le n-1$, $a,k,c,g,h>0$, and
$\Domain=(-1,\infty)^n$. All indices and equalities of indices are
modulo $n$. Write $\mu=a+c-g$ and $A=k+h$ as in the project notation.
P01 supplies smoothness on $\Domain$; P02 supplies the synchronous
restriction and, only where explicitly invoked below, its equilibrium
relation. No formula is imported from the archived blueprint.
The stability equivalences below concern the \emph{linearized system}.
Nonlinear behavior at a zero spectral abscissa and bifurcation type
are reserved for P09. The positive-state Fourier spectrum and the
homogeneous normality theorem are reserved for P04.

\subsection{General Jacobian}

Let $\delta^{(n)}_{j\ell}$ be one if $j\equiv\ell\pmod n$ and
zero otherwise. These are modular Kronecker symbols, not mutually
exclusive case labels.

\begin{lemma}[Jacobian entries]
\label{lem:P03-jacobian}
For every $x\in\Domain$ and $i,j\in\Zn$, $J(x)=DF(x)$ has entries
\begin{align}
 J_{ij}(x)={}&\left(a-2kx_i-\frac{c x_{i-1}}{(1+x_i)^2}
             -\frac{g}{1+x_{i+q}}-h x_{i-q}\right)
                 \delta^{(n)}_{ji}\notag\\
 &+\frac{c}{1+x_i}\delta^{(n)}_{j,i-1}
   +\frac{g x_i}{(1+x_{i+q})^2}\delta^{(n)}_{j,i+q}
   -h x_i\delta^{(n)}_{j,i-q}.
 \label{eq:P03-general-entries}
\end{align}
All contributions are added when indices coincide; all other entries
are zero. The three off-diagonal indices differ from $i$, but need
not differ from each other.
\end{lemma}
\begin{proof}
Differentiate the five terms of \eqref{eq:canonical-dm} by the product
and chain rules. In particular the denominator derivative of
$-gx_i/(1+x_{i+q})$ is $+gx_i/(1+x_{i+q})^2$.
When indices coincide, differentiation adds every contribution,
as encoded by the modular Kronecker symbols.
\end{proof}

For a second formulation, put
\[
 D_i=a-2kx_i-\frac{c x_{i-1}}{(1+x_i)^2}
            -\frac{g}{1+x_{i+q}}-h x_{i-q},\quad
 C_i=\frac{c}{1+x_i},\quad
 G_i=\frac{g x_i}{(1+x_{i+q})^2},\quad H_i=h x_i.
\]
The directional derivative is
\begin{equation}
 (DF(x)z)_i=D_i z_i+C_i z_{i-1}+G_i z_{i+q}-H_i z_{i-q}.
 \label{eq:P03-directional}
\end{equation}
Since $(P^r z)_i=z_{i-r}$ for every integer $r$, this gives
\begin{equation}
 J(x)=\diag(D_i)+\diag(C_i)P+\diag(G_i)P^{-q}
                           -\diag(H_i)P^q.
 \label{eq:P03-variable-matrix}
\end{equation}
The diagonal factors are on the left because their coefficients
belong to receiving row $i$. This general matrix need not be circulant.

\subsection{Linearization at the origin}
\label{slot:P03-1}

Direct substitution gives $F(0)=0$. In
\eqref{eq:P03-general-entries}, $D_i(0)=a-g$, $C_i(0)=c$, and
$G_i(0)=H_i(0)=0$. Both derivatives of the bilinear $h$ term
contain a state coordinate and vanish. The denominator-coordinate
(target) derivative of the input contains $x_{i-1}$ and vanishes;
that of the rational loss contains $x_i$ and vanishes as well.
The numerator derivative of the loss survives as $-g$ on the
 diagonal. Thus, including every collision case,
\begin{equation}
 \boxed{J(0)=(a-g)I+cP.}
 \label{eq:P03-origin-jacobian}
\end{equation}

\begin{proposition}[Origin spectrum]
\label{prop:P03-origin-spectrum}
With the project's DFT convention $\theta_m=2\pi m/n$ and
$v^{(m)}_j=n^{-1/2}e^{\ii\theta_m j}$, the complete spectrum is
\begin{equation}
 \boxed{\lambda_m(0)=a-g+c e^{-\ii\theta_m},\qquad 0\le m<n.}
 \label{eq:P03-origin-spectrum}
\end{equation}
Every eigenvalue is algebraically and geometrically simple over $\C$.
Its real and imaginary parts are
\begin{equation}
 \Ree\lambda_m(0)=\mu-c(1-\cos\theta_m),\qquad
 \Imm\lambda_m(0)=-c\sin\theta_m.
 \label{eq:P03-origin-parts}
\end{equation}
The unique mode attaining the largest real part is $m=0$, and
\begin{equation}
 \boxed{\sa(J(0))=\lambda_0(0)=\mu=a+c-g.}
 \label{eq:P03-abscissa}
\end{equation}
\end{proposition}
\begin{proof}
Apply Gray's circulant diagonalization theorem
\cite[Theorem~3.1]{Gray2006} to \eqref{eq:P03-origin-jacobian}.
With our positive-exponent columns, $Pv^{(m)}=e^{-\ii\theta_m}v^{(m)}$;
Gray's opposite exponent is reconciled by relabelling $m\mapsto -m$.
Distinct roots of unity and $c>0$ give simple eigenvalues. In particular
\begin{equation}
 \det(\lambda I-J(0))=(\lambda-(a-g))^n-c^n.
 \label{eq:P03-characteristic}
\end{equation}
Taking real and imaginary parts gives the stated formulas. Since
$1-\cos\theta_m>0$ for every $m\ne0$, only $m=0$ attains $\mu$.
\end{proof}

\subsection{Origin stability and the nonhyperbolic boundary}
\label{slot:P03-2}

\begin{proposition}[Exact linear stability and hyperbolicity conditions]
\label{prop:P03-stability}
For $\dot z=J(0)z$, zero is asymptotically (indeed exponentially)
stable if and only if $\mu<0$, and unstable if and only if $\mu>0$.
At $\mu=0$ it is Lyapunov stable but not asymptotically stable.
The origin of the nonlinear field is nonhyperbolic if and only if
\begin{equation}
 g=a+c\cos\frac{2\pi m}{n}
 \quad\hbox{for some }m\in\{0,\ldots,n-1\},
 \label{eq:P03-all-nonhyperbolic}
\end{equation}
with $g>0$ as always. Nonhyperbolicity is therefore not restricted
to the stability boundary $\mu=0$: every locus in
\eqref{eq:P03-all-nonhyperbolic} with $m\ne0$ has $\mu>0$.
\end{proposition}
\begin{proof}
In the orthonormal basis of Proposition~\ref{prop:P03-origin-spectrum},
each coefficient solves $\dot{\widehat z}_m=\lambda_m(0)\widehat z_m$.
For $\mu<0$, all coefficients decay exponentially, with
$\norm{z(t)}_2\le e^{\mu t}\norm{z(0)}_2$ for $t\ge0$.
For $\mu>0$, the real synchronous solution
$z(t)=\varepsilon e^{\mu t}\one$ proves instability for arbitrarily
small $\varepsilon>0$. For $\mu=0$, every nonsynchronous coefficient
 decays and the synchronous coefficient is constant. Therefore
$\norm{z(t)}_2\le\norm{z(0)}_2$, while nonzero constant synchronous
solutions prevent attraction. This also proves both converses
without a nonlinear stability theorem. The definition of hyperbolicity
and \eqref{eq:P03-origin-parts} give
\eqref{eq:P03-all-nonhyperbolic}. For $m\ne0$ its value of
$\mu$ is $c(1-\cos\theta_m)>0$.
\end{proof}

The spectral-abscissa sign and the linear stability threshold depend
on neither $n$, $q$, nor $\gcd(n,q)$. Individual eigenvalues and the
\emph{full set} of nonhyperbolic loci depend on $n$, but not on $q$
or a coprimality condition. For example, $n=6$, $a=2$, $c=1$,
$g=5/2$ gives $\mu=1/2>0$ and purely imaginary eigenvalues in
modes $1,5$; mode zero is already unstable. This is not a Hopf claim.
More generally, using $0\le m\le\lfloor n/2\rfloor$ as pair
representatives, each interior representative in
\eqref{eq:P03-all-nonhyperbolic} contributes exactly the two
simple nonzero imaginary eigenvalues in modes $m,n-m$.
Strict decrease of the cosine on $[0,\pi]$ excludes any other
simultaneous zero-real-part mode. The endpoint cases are treated below.

\begin{proposition}[Simple zero at the stability boundary]
\label{prop:P03-boundary}
At $g=a+c$ the only zero eigenvalue is $\lambda_0(0)=0$.
Its algebraic and geometric multiplicities are both one, with
kernel $\operatorname{span}\{\one\}$. Every nonzero mode has
strictly negative real part; none can simultaneously have zero real part.
\end{proposition}
\begin{proof}
At $g=a+c$, \eqref{eq:P03-origin-spectrum} becomes
$c(e^{-\ii\theta_m}-1)$. Only $m=0$ is zero, and every other real
part is negative. Simplicity follows from Proposition~\ref{prop:P03-origin-spectrum};
$Pz=z$ is equivalent to $z\in\operatorname{span}\{\one\}$.
\end{proof}

\begin{remark}[A zero eigenvalue away from the stability boundary]
\label{rem:P03-other-zero}
Over \emph{all} canonical parameter values, a nonzero-mode zero
is possible. Its imaginary part requires $\sin\theta_m=0$;
the only nonzero possibility is even $n$ and $m=n/2$.
The real part then vanishes exactly at $g=a-c$, which is
canonical precisely when $a>c$. That zero is also simple, but
$\mu=2c>0$, so the origin is already linearly unstable. For example,
$a=2$, $c=g=1$ works in every even canonical dimension.
Thus $c>0$ excludes nonzero-mode zeros at $g=a+c$, not at every
parameter value. For odd $n$ there is no nonzero-mode zero.
\end{remark}

These are spectral and linear-system conclusions. ``Nonhyperbolic''
and ``linearly unstable'' are not disjoint categories. The Hurwitz
criterion is not an assertion that nonlinear asymptotic stability
holds if and only if $\mu<0$. Linearization alone does not decide
nonlinear behavior at $\mu=0$, either in $\Domain$ or relative to
$\Orth$. The simple zero does not establish a transcritical
bifurcation or a nonlinear nondegeneracy condition.

\subsection{Jacobian at a synchronous state}
\label{slot:P03-3}
Let $x=x^*\one$ be any synchronous state with $x^*>-1$ and set
$s=1+x^*>0$. Here $s$ denotes the positive denominator, not the
scalar state variable denoted by $s$ in P02.

\begin{proposition}[Homogeneous Jacobian]
\label{prop:P03-homogeneous}
At every such state, whether or not it is an equilibrium, define
\begin{align}
 d&=a-2kx^*-\frac{c x^*}{s^2}-\frac{g}{s}-h x^*,
       \label{eq:P03-d}\\
 C&=\frac{c}{s},\qquad G=\frac{g x^*}{s^2},\qquad H=h x^*.
       \label{eq:P03-CGH}
\end{align}
Then
\begin{equation}
 \boxed{J(x^*\one)=dI+CP+GP^{-q}-HP^q.}
 \label{eq:P03-homogeneous}
\end{equation}
Its entries are
$d\delta^{(n)}_{ji}+C\delta^{(n)}_{j,i-1}
+G\delta^{(n)}_{j,i+q}-H\delta^{(n)}_{j,i-q}$,
with sums retained at coincident indices.
\end{proposition}
\begin{proof}
Substitute $x_i=x^*$ in \eqref{eq:P03-general-entries}.
The index $i+q$ corresponds to $P^{-q}$ and $i-q$ to $P^q$;
coincident terms still add.
\end{proof}

Always $C>0$. If $x^*>0$, then $G,H>0$, and the separate
neighbor contributions have signs $+C,+G,-H$. If $x^*=0$,
then $G=H=0$; if $-1<x^*<0$, then $G,H<0$.
No fixed sign for $d$ at an arbitrary synchronous state is assumed.
Signs at colliding offsets must be assessed after addition.

\paragraph{Optional equilibrium simplification using P02.}
Only now suppose $x^*\ne0$ and $x^*\one$ is an equilibrium.
P02's \eqref{eq:P02-scalar}, divided by this nonzero state, gives
\begin{equation}
 a-(k+h)x^*+\frac{c-g}{1+x^*}=0.
 \label{eq:P03-equilibrium-relation}
\end{equation}
Substitution of $a=(k+h)x^*-(c-g)/s$ in the unsimplified
 diagonal coefficient yields
\begin{equation}
 d=-k x^*-\frac{c}{s}-\frac{c x^*}{s^2}
   =-k x^*-\frac{c(1+2x^*)}{s^2}.
 \label{eq:P03-equilibrium-d}
\end{equation}
This applies in particular at a positive synchronous equilibrium.
It is not available at an arbitrary state or at the origin merely
because the origin is an equilibrium. No stability conclusion is drawn.

\subsection{Circulant representation and offset collisions}

Since $P^n=I$, \eqref{eq:P03-homogeneous} is the ordinary polynomial
$dI+CP+GP^{n-q}-HP^q$. Each power of $P$ has entries depending
only on the cyclic difference $i-j$; their linear combination has
rows obtained by cyclic shifts of the preceding row. This proves
that the homogeneous Jacobian is circulant without requiring four
distinct cyclic diagonals. The standard consequence is Fourier
 diagonalizability; the actual homogeneous spectrum and normality
 theorem are left to P04.

The diagonal cannot collide with a neighbor since $1\not\equiv0$
and $q\not\equiv0$ modulo $n$. Neighbor collisions are exhausted by
\[
\begin{aligned}
 i-1\equiv i-q&\iff q=1,\\
 i-1\equiv i+q&\iff q=n-1,\\
 i+q\equiv i-q&\iff 2q\equiv0\pmod n.
\end{aligned}
\]
The last case means even $n$ and $q=n/2$. For $n\ge5$, these
three cases are mutually exclusive; no triple collision occurs.
A nontrivial $\gcd(n,q)$ creates no further coincident offsets.
The resulting matrices are
\begin{align}
 q=1:&\quad J=dI+(C-H)P+GP^{-1},
    \label{eq:P03-collision-one}\\
 q=n-1:&\quad J=dI+(C+G)P-HP^{-1},
    \label{eq:P03-collision-last}\\
 n\text{ even},\ q=n/2:&\quad J=dI+CP+(G-H)P^{n/2}.
    \label{eq:P03-collision-half}
\end{align}
For other canonical $q$ the three neighbor offsets are distinct.
Even then a coefficient may vanish; the representation lists
contributions, not a claim that every entry is nonzero.

At nonsynchronous states the same rule gives, for $q=1$, the
entry $c/(1+x_i)-h x_i$ at $j=i-1$; for $q=n-1$, the entry
$c/(1+x_i)+g x_i/(1+x_{i-1})^2$ at $j=i-1$; and for even $n$,
$q=n/2$, the entry $g x_i/(1+x_{i+q})^2-h x_i$ at
$j=i+q=i-q$. The actual repeated coordinate is identified before
combining contributions. At homogeneous states $C-H$ or $G-H$
can vanish through cancellation. The summed polynomial remains valid.

Differentiating P02's synchronous field also gives the consistency identity
\begin{equation}
 f'(r)=a-2(k+h)r+\frac{c-g}{(1+r)^2}
       =d+C+G-H\quad(x^*=r).
 \label{eq:P03-row-sum-check}
\end{equation}
